% !PRINTED 313
\weberpage{313}
%<<CONT
auf die ersten Fälle beschränken, und wir beginnen also mit den quadratischen Gleichungen
\weq{1}{x^2 + a\,x + b = 0.}
Sehen wir $a$ und $b$ als rechtwinklige Coordinaten in einer Ebene an, so ist jeder Punkt dieser Ebene das Bild einer und nur einer quadratischen Gleichung, und umgekehrt hat jede quadratische Gleichung mit reellen Co\"efficienten einen Punkt der Ebene zum Bilde.

Die Gleichung
\weq{2}{a^2 - 4\,b = 0}
stellt eine Parabel dar, auf der die Bilder der Gleichungen mit gleichen Wurzeln liegen. Die inneren Punkte, d. h. die in der H\"ohlung der Parabel liegenden, stellen die Gleichungen mit imaginären Wurzeln, die äusseren Punkte die Gleichungen mit reellen Wurzeln dar.

Wir bilden nun nach \S.\,89 die \gesperrt{Sturm'sche} Kette
\weq{3}{x^2 + a\,x + b,\quad 2\,x + a,\quad a^2 - 4\,b}
und untersuchen die Zeichenwechsel. Ist $x$ constant und $a$, $b$ variabel, so ist
\[
x^2 + a\,x + b = 0
\]
die Gleichung der Parabeltangente in dem Punkte
\[
a = -2\,x,\quad b = x^2,
\]
und zwar ist $x^2 + a\,x + b$ positiv in dem Theil der Ebene, der zugleich die Parabel enthält.

Durch diese Tangente und durch die Parabel wird die Ebene in vier Felder getheilt, in denen in der Reihe (3), wie in der Fig. 16 (a. f. S.) angedeutet ist, 0, 1, 2 Zeichenwechsel vorkommen.

Wenn man also $x = \alpha$ und $x = \beta$ setzt, so erhält man zwei solche Tangenten und eine Eintheilung der Ebene in vier Felder, worin die repräsentirenden Punkte für solche Gleichungen liegen, die zwischen $\alpha$ und $\beta$ keine, eine oder zwei Wurzeln haben. In der Fig. 17 sind diese vier Felder durch die betreffenden Ziffern gekennzeichnet.

Es ist der Unterschied zwischen solchen Punkten, von denen aus keine, eine oder zwei Tangenten an die Parabel gezogen werden können, die ihren Berührungspunkt zwischen $\alpha$ und $\beta$ haben. Die Fig. 17 also giebt die genaue und vollständige Unterscheidung der verschiedenen Fälle.
